\documentclass[12pt]{article}
\usepackage[margin=1in]{geometry}
\usepackage{parskip}
\pdfinfoomitdate=1
\pdftrailerid{}
\begin{document}

Consider two certificates $P$ and $Q$ where $P$ is signed by $Q$.
Let $s$ be the signature in $P$ and let $k$ be the public
key in $Q$.

\begin{center}
\begin{tabular}{|c|}
\hline
$\quad P\quad$\\
\\
\\
\\
\\
\\
\hline
$s$\\
\hline
\end{tabular}
\qquad
\begin{tabular}{|c|}
\hline
$\quad Q\quad$\\
\\
\\
\hline
$k$\\
\hline
\\
\\
\\
\hline
\end{tabular}
\end{center}

Let $k'$ be the secret key associated with $k$.
Only the owner of $Q$ knows the secret key.

Signature $s$ is a hash digest of $P$ encrypted using secret key $k'$.
For example, here is a 74 byte signature in ASN.1 notation for prime256v1
and hash digest sha256.

{\scriptsize
\begin{verbatim}
233: tag 0x03,   72 bytes BIT STRING 00 
236: tag 0x30,   69 bytes SEQUENCE
238: tag 0x02,   32 bytes INTEGER 792586f140a100dd18b2189a00d774d55105907c5ae4a67069ead8d871f17f85
272: tag 0x02,   33 bytes INTEGER 00ff3dab62841cc33530454149fa05cca5bccc80d389ea9657b197f6d763954a9d
\end{verbatim}
}

The two integers are the encryption result.
The second integer is 33 bytes because of the ASN.1 coding rule that the first bit
of an integer must be 0.
Note the ASN.1 object hierarchy of ``sequence'' encompassing the two integers
and ``bit string'' encompassing the ``sequence'' object.

To prove that $P$ is signed by $Q$, signature $s$ is decrypted using $Q$'s public key $k$.
Then if the decrypted value matches the digest of $P$ (sha256 in this example),
the signing of $P$ by $Q$ is proven.
The contents of $P$ cannot be changed without breaking signature $s$,
and $s$ cannot be changed without knowing the secret key $k'$.
Hence the contents of $P$ can be trusted.

\end{document}
